% ============================================================================
% Fichier  : partie4/P04_C15_audit_projets.tex
% Titre    : Auditer les Projets Informatiques
% Auteur   : MINKA MI NGUIDJOÏ Thierry Emmanuel
% Version  : 1.0 -- novembre 2026
% Base   : ACSI C15, repris et développé pour INF3226
% ============================================================================
% Sources : guide d'audit des SI du CHAI (2014), fiches 3.6 « fonction
% étude », 3.7 « projets » (objectifs et enjeux, opportunité et besoins,
% planification, instances de pilotage, méthodes, plan d'assurance qualité,
% tests et recette, gestion des évolutions, mise en production) et 3.8
% « marchés » (risques des marchés informatiques), paraphrasées et
% transposées hors du droit français ; COBIT 2019 (APO05, BAI01, BAI03,
% BAI07, BAI11) comme critères. Faits HORIZON (BANQUE SANAGA) et ANPP repris
% du cours GCI. Méthode de la valeur acquise : technique classique de
% gestion de projet, présentée sans référence normative.
% ============================================================================

\chapter{Auditer les Projets Informatiques}
\label{chap:audit-projets}

\epigraph{%
    « Un projet ne dérape pas le jour de la bascule.\\
    Il dérape le jour où l'on cesse de le regarder. »%
}{Adage des chefs de projet}

\niveauChapitre{Intermédiaire}{%
    Maîtrise du Chapitre~\ref{chap:audit-gouvernance} (gouvernance,
    organes, indicateurs) et du Chapitre~\ref{chap:audit-applications}
    (contrôles d'application)%
}

\begin{boxobjectifs}
\begin{itemize}
    \item Distinguer l'audit d'un projet, l'audit de la fonction « études »
          et l'audit d'une application en service.
    \item Décomposer l'audit d'un projet, des enjeux à la mise en
          production.
    \item Auditer la décision de mise en production~: critères, organe,
          retour arrière.
    \item Identifier les signaux d'un projet en dérive, et les mesurer par
          la valeur acquise.
    \item Situer les points de contrôle de la fonction « études » et des
          marchés informatiques.
\end{itemize}
\end{boxobjectifs}

\bigskip

% ============================================================================
\section*{Question directrice : quand fallait-il regarder~?}
% ============================================================================

À la BANQUE SANAGA, le projet HORIZON a remplacé le système bancaire
central pour 3,2~milliards FCFA. La bascule du 7~mars~2026 a privé les
agences de service pendant trois jours (cours MSN-2S3-MSI03). Le comité
d'audit demande à Mme~EKAMBI~: « Qu'aurait vu un audit du projet, et
quand~? »

À l'ANPP, le schéma directeur de 2023 prévoyait neuf projets pour
14~milliards FCFA. À mi-parcours, trois projets ont consommé 71~\% du
budget sans être achevés. La Chambre des comptes veut savoir s'il est
encore temps d'agir.

\medskip
\textbf{Plan du chapitre.}
\begin{enumerate}
    \item Décomposer l'audit d'un projet (section~\ref{sec:c19-decomposer}).
    \item Auditer la mise en production (section~\ref{sec:c19-production}).
    \item Repérer et mesurer la dérive (section~\ref{sec:c19-derive}).
    \item La fonction « études » et les marchés (section~\ref{sec:c19-etudes}).
    \item Auditer un marché informatique, étape par étape (section~\ref{sec:c19-marche}).
\end{enumerate}

\separateur

% ============================================================================
\section{Décomposer l'Audit d'un Projet}
\label{sec:c19-decomposer}
% ============================================================================

\subsection{Projet, fonction « études », application}

Le guide du CHAI insiste sur une distinction~: l'audit d'un \emph{projet}
porte sur une organisation temporaire et sur mesure, créée pour atteindre
un objectif dans un délai~; l'audit de la \emph{fonction « études »} porte
sur l'organisation permanente qui développe et maintient les
applications~\cite{chai2014guide}. L'audit d'une \emph{application en
service} (Chapitre~\ref{chap:audit-applications}) commence quand le projet
s'achève.

Un projet peut être audité à deux moments~:
\begin{itemize}
    \item \textbf{en cours}, aux jalons~: l'audit peut encore changer le
          cours du projet~;
    \item \textbf{a posteriori}~: il établit ce qui s'est passé, et en tire
          des leçons pour les projets suivants.
\end{itemize}

\subsection{Cinq sous-domaines}

\begin{tableaucours}{cL{3.0cm}L{5.8cm}L{2.0cm}}{%
    \tableauHeader{Code} &
    \tableauHeader{Sous-domaine} &
    \tableauHeader{Contrôles (d'après le guide du CHAI)} &
    \tableauHeader{COBIT}
}
PR.1 & Enjeux et besoins & Étude d'opportunité et de valeur~; scénarios
comparés, y compris celui de ne rien faire~; cahier des charges rédigé par
la maîtrise d'ouvrage & APO05, BAI03 \\
PR.2 & Planification et pilotage & Plan de projet révisé~; analyse des
risques du projet~; comité de pilotage, comité de projet, comité des
utilisateurs actifs, avec leurs délégations & BAI01, BAI11 \\
PR.3 & Qualité et méthode & Méthode appliquée~; plan d'assurance qualité~;
groupe qualité indépendant de l'équipe~; évolutions du périmètre
formalisées et chiffrées & BAI11 \\
PR.4 & Tests et recette & Plan de tests~; recette par la maîtrise
d'ouvrage~; procès-verbal~; reprise des données testée~; tests de
pré-exploitation & BAI07 \\
PR.5 & Mise en production & Critères de passage écrits~; décision par
l'organe compétent~; retour arrière préparé et testé~; transfert de
responsabilité à la production & BAI07 \\
\end{tableaucours}
\vspace{-0.8em}
\begin{center}{\small\itshape Tableau~19.1. Décomposition de l'audit d'un
projet}\end{center}

Pour chaque contrôle, le guide du CHAI énumère les documents à récupérer~:
étude de valeur, plan de projet, comptes rendus des comités, plan
d'assurance qualité, plans et comptes rendus de tests et de recette,
dossier de reprise des données, procédure de mise en production. Leur
absence est déjà une information~: un comité sans compte rendu est, pour
l'auditeur, un comité qui n'a rien décidé.

% ============================================================================
\section{Auditer la Mise en Production}
\label{sec:c19-production}
% ============================================================================

\subsection{Une décision, pas une date}

La mise en production est le moment où le risque du projet devient le
risque de l'organisation. Elle doit résulter d'une décision, prise par
l'organe compétent, sur des critères écrits à l'avance~: taux de scénarios
recettés, anomalies bloquantes résolues, reprise des données validée,
retour arrière testé, équipes de production formées.

\begin{boxprocedure}
\textbf{Auditer une décision de mise en production.}
\begin{enumerate}
    \item \textbf{Critères}~: une grille de passage existe, approuvée avant
          la date de décision.
    \item \textbf{Mesure}~: chaque critère est mesuré, pièce à l'appui
          (procès-verbal de recette, rapport de reprise, test de retour
          arrière).
    \item \textbf{Organe}~: la décision est prise par le comité de pilotage,
          et non par une personne seule.
    \item \textbf{Écarts}~: un critère non atteint est accepté par écrit,
          avec le risque nommé et la mesure de compensation.
    \item \textbf{Retour arrière}~: un point de non-retour est défini, et la
          procédure de retour a été testée avant ce point.
\end{enumerate}
\end{boxprocedure}

\begin{boxlecture}{Projet HORIZON, a posteriori}
Mme~EKAMBI cote a posteriori les vingt contrôles du projet HORIZON, d'après
la chronologie du cours GCI~: recette à 63~\% des scénarios le
20~février~; dernière réunion du comité de pilotage en novembre~2025~;
bascule fixée par le directeur général pour coïncider avec un lancement
commercial~; note de report de la DSI restée sans réponse écrite~;
écarts de solde sur 11~\% des comptes, sans retour arrière possible.

\medskip
\begin{tableaucours}{L{4.2cm}ccc}{%
    \tableauHeader{Sous-domaine} &
    \tableauHeader{\OUI{}} & \tableauHeader{\NON{}} & \tableauHeader{\NEK{}}
}
PR.1 Enjeux et besoins & 3 & 1 & 0 \\
PR.2 Planification et pilotage & 1 & 3 & 0 \\
PR.3 Qualité et méthode & 2 & 1 & 1 \\
PR.4 Tests et recette & 1 & 2 & 1 \\
PR.5 Mise en production & 0 & 4 & 0 \\
Projet HORIZON & 7 & 11 & 2 \\
\end{tableaucours}

\medskip
\textbf{Constatation.} $\mu = 0{,}35$~; $U = 0{,}10$~; $H \approx 1{,}34$~bit.
Les quatre contrôles de PR.5 sont \NON{}. Les deux \NEK{} portent sur
l'indépendance du contrôle qualité et sur les tests de reprise des
données, dont les comptes rendus n'ont pas été retrouvés.

\medskip
\textbf{Analyse} (confiance élevée). Non-conformité claire, concentrée dans
le pilotage et la mise en production. Le projet n'a pas manqué de
méthode~: il a manqué d'un organe qui décide. Un audit au jalon de février
aurait relevé un comité de pilotage inactif depuis trois mois et une
recette incomplète~: deux faits suffisants pour demander une décision
écrite du comité.

\medskip
\hfill\lecture{0{,}35}{0{,}10}{1{,}34}
\end{boxlecture}

\begin{boxobjet}{PR.5.2~: Observation}
\textbf{1. Norme~:} procédure projets de la banque~: « la mise en
production est décidée par le comité de pilotage, sur la grille de
passage »~; COBIT BAI07.

\textbf{2. Constatation~:} la date de bascule a été fixée en janvier~2026
par courriel du directeur général. Le comité de pilotage ne s'est pas
réuni après novembre~2025. Aucune grille de passage n'a été produite.
Cote \NON{}.

\textbf{3. Risques~:} mise en production d'un système non prêt~; perte de
service~; atteinte à l'intégrité des soldes.

\textbf{4. Contradictoire~:} la direction générale indique, le 14~avril,
que la date était «~imposée par le marché~».

\textbf{5. Analyse} (confiance élevée)~: la contrainte commerciale a
remplacé la décision collégiale~; la note de report de la DSI n'a reçu
aucune réponse écrite.

\textbf{6. Conclusion~:} le contrôle de la décision de mise en production
n'a pas fonctionné.

\textbf{7. Recommandations~:} grille de passage obligatoire, approuvée par
le comité de pilotage~; toute dérogation acceptée par écrit avec le risque
nommé~; information du comité d'audit avant toute bascule d'un système
critique.
\end{boxobjet}

\begin{boxalert}
\textbf{Violation typique.} L'auditeur, présent au comité de pilotage,
déclare que « le projet peut basculer ».

\smallskip
\textbf{Règle~:} l'auditeur dit si la décision est prise selon les règles
et sur des critères mesurés~; il ne prend pas la décision. Sinon, il
auditera plus tard sa propre décision~: c'est une menace
d'autorévision (Chapitre~\ref{chap:normes}).
\end{boxalert}

% ============================================================================
\section{Repérer et Mesurer la Dérive}
\label{sec:c19-derive}
% ============================================================================

\subsection{Les signaux}

\begin{tableaucours}{L{4.6cm}L{4.0cm}L{3.6cm}}{%
    \tableauHeader{Signal} &
    \tableauHeader{Preuve} &
    \tableauHeader{Dans nos cas}
}
Le comité de pilotage ne se réunit plus & Calendrier et comptes rendus &
HORIZON~: rien après novembre~2025 \\
La date est fixée par un événement extérieur au projet & Courriels,
annonces & HORIZON~: lancement commercial \\
La recette est réduite pour tenir la date & Plan de tests, procès-verbal &
HORIZON~: 63~\% \\
Une alerte écrite reste sans réponse & Notes internes & HORIZON~: note de
la DSI \\
Le budget consommé dépasse l'avancement & Suivi financier, avancement &
ANPP~: 71~\% du budget pour trois projets \\
Le périmètre change sans réestimation & Demandes d'évolution & À vérifier
à l'ANPP \\
Le chef de projet change plusieurs fois & Organigrammes du projet & À
vérifier à l'ANPP \\
\end{tableaucours}
\vspace{-0.8em}
\begin{center}{\small\itshape Tableau~19.2. Signaux d'un projet en
dérive}\end{center}

\subsection{La valeur acquise}

Le signal financier se mesure. La méthode de la valeur acquise, classique
en gestion de projet, compare trois grandeurs~: le budget prévu à
l'achèvement (\emph{BAC}), le coût réel déjà consommé (\emph{CR}) et la
\emph{valeur acquise} (\emph{VA}), c'est-à-dire la part du budget qui
correspond au travail réellement accompli~:
\[
    VA = \text{avancement} \times BAC, \qquad
    IPC = \frac{VA}{CR}, \qquad
    EAC = \frac{BAC}{IPC}.
\]
Un indice de performance des coûts (\emph{IPC}) inférieur à 1 signifie
que chaque franc dépensé produit moins d'un franc de travail prévu~;
\emph{EAC} estime le coût final si cette performance se maintient.

\begin{boxscenario}{Opération MEFOU. Trois projets de l'ANPP}
\textbf{Faits} (suivi financier et rapports d'avancement validés par la
maîtrise d'ouvrage, montants en milliards FCFA).

\medskip
\begin{tableaucours}{L{3.6cm}cccccc}{%
    \tableauHeader{Projet} &
    \tableauHeader{BAC} & \tableauHeader{Avanc.} & \tableauHeader{VA} &
    \tableauHeader{CR} & \tableauHeader{IPC} & \tableauHeader{EAC}
}
Paie et pensions & 3,50 & 60~\% & 2,10 & 3,90 & 0,54 & 6,50 \\
Dématérialisation & 2,60 & 45~\% & 1,17 & 2,84 & 0,41 & 6,31 \\
Réseau et hébergement & 3,20 & 70~\% & 2,24 & 3,20 & 0,70 & 4,57 \\
Total & 9,30 & -- & 5,51 & 9,94 & 0,55 & 17,38 \\
\end{tableaucours}

\medskip
\textbf{Constatation.} Les trois projets ont consommé 9,94~milliards,
soit 71~\% des 14~milliards du schéma directeur, pour une valeur acquise
de 5,51~milliards. À performance constante, ils coûteraient environ
17,4~milliards, près du double de leur budget.

\medskip
\textbf{Analyse} (confiance moyenne~: les taux d'avancement sont déclarés
par la maîtrise d'ouvrage et doivent être vérifiés sur les livrables). Le
schéma directeur ne peut plus être tenu dans son enveloppe~: les six
autres projets ne disposent plus que de 4,06~milliards. La décision
d'arbitrer appartient au conseil d'administration~; l'auditeur doit la
lui rendre possible, chiffres à l'appui.
\end{boxscenario}

\begin{boxinfo}
\textbf{L'avancement est une affirmation.} La valeur acquise repose sur un
taux d'avancement déclaré. Comme tout indicateur
(Chapitre~\ref{chap:audit-gouvernance}), il se vérifie~: un avancement
non adossé à des livrables recettés se cote \NEK{}, et l'auditeur présente
alors une fourchette plutôt qu'un chiffre.
\end{boxinfo}

% ============================================================================
\section{La Fonction « Études » et les Marchés}
\label{sec:c19-etudes}
% ============================================================================

\subsection{La fonction « études »}

La fonction « études » développe et maintient les applications. Le guide
du CHAI en propose un audit de la capacité à maintenir le patrimoine
applicatif et à conduire les évolutions du schéma
directeur~\cite{chai2014guide}. Les points essentiels pour l'auditeur des
SI sont les suivants~:
\begin{itemize}
    \item une ligne de partage écrite entre études et production, et des
          environnements de développement, de test et de production
          séparés~;
    \item une procédure de maintenance~: demande validée par le
          propriétaire de l'application, estimation, priorisation, recette~;
    \item des évolutions regroupées en versions, plutôt qu'au fil de
          l'eau, pour limiter le risque de régression~;
    \item des corrections urgentes tracées et revues après coup~;
    \item une gestion des versions du code, qui permet de savoir qui a
          modifié quoi.
\end{itemize}

\subsection{Les marchés informatiques}

Les projets sont souvent réalisés par un prestataire, dans le cadre d'un
marché public ou d'un contrat. Le guide du CHAI relève des risques propres
à ces marchés~: responsabilités floues entre maîtrise d'ouvrage,
maîtrise d'œuvre et assistance, obligations de moyens là où l'on attendait
des résultats, périmètre imprécis qui enferme l'acheteur dans une
dépendance. Au Cameroun, le cadre est le Code des marchés publics (texte
en vigueur à vérifier au Journal officiel)~; les points de contrôle
contractuels restent les mêmes.

\begin{tableaucours}{L{4.4cm}L{7.8cm}}{%
    \tableauHeader{Point de contrôle} &
    \tableauHeader{Ce que l'auditeur vérifie}
}
Objet et responsabilités & Les rôles de l'acheteur et du titulaire sont
écrits et connus des équipes \\
Livrables et jalons & Le marché exige des résultats mesurables, et non un
nombre de jours de présence \\
Réception & Chaque livrable est réceptionné par procès-verbal, avec les
réserves \\
Évolutions & Toute modification du périmètre fait l'objet d'un avenant
chiffré \\
Réversibilité & Le contrat prévoit la restitution des données, de la
documentation et, pour un développement spécifique, des sources ou leur
dépôt chez un tiers \\
Garantie et pénalités & Les pénalités sont prévues et, en cas de retard,
effectivement appliquées \\
\end{tableaucours}
\vspace{-0.8em}
\begin{center}{\small\itshape Tableau~19.3. Points de contrôle d'un marché
informatique (d'après le guide du CHAI)}\end{center}

\begin{boxmetier}{Auditeur SI}
\textbf{Auditer un projet à un jalon}
\begin{enumerate}
    \item Demander le plan de projet en vigueur, les trois derniers comptes
          rendus du comité de pilotage et le tableau de bord du projet.
    \item Vérifier que chaque comité prévu s'est réuni, et que ses comptes
          rendus contiennent des décisions.
    \item Recalculer la valeur acquise à partir des livrables recettés, et
          non des taux déclarés.
    \item Lire la grille de passage du jalon suivant~; s'il n'y en a pas,
          le signaler avant la date.
    \item Rendre un avis écrit sur le processus de décision, sans prendre
          la décision.
\end{enumerate}
\end{boxmetier}

\begin{boxmetier}{Gestionnaire de risques}
\textbf{Suivre un portefeuille de projets}
\begin{enumerate}
    \item Tenir pour chaque projet le budget, le coût réel et la valeur
          acquise, mis à jour chaque trimestre.
    \item Fixer un seuil d'alerte, par exemple un IPC inférieur à 0,8.
    \item Porter tout projet en alerte au registre des risques de
          l'organisation, coté sur l'échelle commune.
    \item Soumettre au conseil, en cas d'alerte, les options~: poursuivre,
          réduire le périmètre, suspendre.
\end{enumerate}
\end{boxmetier}

% ============================================================================

% ============================================================================
\section{Auditer un Marché Informatique}
\label{sec:c19-marche}
% ============================================================================

Un marché informatique public suit des étapes que l'auditeur vérifie l'une
après l'autre~; les règles précises sont celles du code des marchés publics
applicable (au Cameroun, texte à vérifier dans sa version en vigueur).

\begin{tableaucours}{L{2.8cm}L{4.6cm}L{4.8cm}}{%
    \tableauHeader{Étape} &
    \tableauHeader{Question d'audit} &
    \tableauHeader{Preuve}
}
Besoin & Le besoin est-il défini avant le choix d'une solution~? & Étude
préalable, expression des besoins datée \\
Cahier des charges & Les exigences sont-elles mesurables, et neutres
vis-à-vis des fournisseurs~? & Cahier des charges, critères de recette \\
Mise en concurrence & La procédure prévue a-t-elle été suivie~? & Avis
publié, registre des offres, procès-verbal d'ouverture \\
Attribution & Les critères annoncés ont-ils été appliqués~? & Rapport
d'analyse des offres, décision d'attribution \\
Exécution & Les livrables sont-ils conformes au contrat, et payés après
réception~? & Procès-verbaux de recette, factures, paiements \\
Réception et réversibilité & Le maître d'ouvrage peut-il exploiter et faire
évoluer le système sans le titulaire~? & Documentation, codes sources ou
séquestre, clause de réversibilité \\
\end{tableaucours}
\vspace{-0.8em}
\begin{center}{\small\itshape Tableau~19.4. Les étapes d'un marché
informatique et leurs preuves}\end{center}

\begin{boxalert}
\textbf{Violation typique.} Un cahier des charges qui cite une marque ou un
produit (« solution basée sur le logiciel X »), puis une mise en
concurrence que seul le revendeur de X peut gagner.

\smallskip
\textbf{Règle~:} les exigences décrivent un besoin fonctionnel et des
performances mesurables~; une référence à un produit, si elle est
inévitable, est suivie de « ou équivalent ». Un critère qui oriente le choix
est un constat sur la mise en concurrence.
\end{boxalert}

\begin{boxscenario}{Opération MEFOU. Le marché de la dématérialisation}
\textbf{Faits.} Le projet de dématérialisation de l'ANPP (2,60~milliards
FCFA, 45~\% d'avancement) a été attribué après un appel d'offres dont le
cahier des charges exige « une compatibilité native » avec le produit d'un
éditeur. Une seule offre a été reçue.

\medskip
\textbf{Analyse} (confiance moyenne). Le critère a pu restreindre la
concurrence~; une offre unique est un signal, pas une preuve. M.~NJOYA
demande le rapport d'analyse et les échanges préalables avec l'éditeur, et
cote \NEK{} le contrôle « la mise en concurrence a été effective » en
attendant.
\end{boxscenario}

\separateur

\begin{boxresume}
\textbf{Ce qu'il faut retenir de ce chapitre~:}
\begin{itemize}
    \item \textbf{Projet, études, application}~: trois objets d'audit~; le
          projet est une organisation temporaire, auditée en cours ou a
          posteriori.
    \item \textbf{Cinq sous-domaines}~: enjeux et besoins, planification et
          pilotage, qualité et méthode, tests et recette, mise en
          production.
    \item \textbf{La mise en production est une décision}~: critères écrits,
          organe compétent, écarts acceptés par écrit, retour arrière testé.
    \item \textbf{L'auditeur ne décide pas} de la bascule~: il dit si la
          décision est prise selon les règles.
    \item \textbf{La dérive se repère} (comités muets, recette réduite,
          alertes sans réponse) \textbf{et se mesure} par la valeur acquise.
    \item \textbf{Les marchés} se contrôlent par les livrables, la réception,
          les avenants, la réversibilité et les pénalités.
    \item \textbf{Un marché informatique} s'audite étape par étape, du besoin à la réversibilité~; un cahier des charges qui cite un produit oriente la concurrence.
\end{itemize}
\end{boxresume}

\bigskip

\begin{boxexo}[title={\ding{45}\quad Exercice~19.1. Rattacher les contrôles \niveaubase}]
Rattachez chaque contrôle à un sous-domaine (PR.1 à PR.5).
\begin{enumerate}[(a)]
    \item Le scénario « ne rien faire » a été chiffré.
    \item La reprise des données a été testée sur une copie complète.
    \item Le comité des utilisateurs valide les dossiers de conception.
    \item Le retour arrière a été répété avant le point de non-retour.
    \item Le groupe qualité ne rend pas compte au chef de projet.
\end{enumerate}
\end{boxexo}

\begin{boxexo}[title={\ding{45}\quad Exercice~19.2. Calculer la valeur acquise \niveaumoyen}]
Un projet de LiZbiZ a un budget de 800~millions FCFA. Au 30~juin, il a
consommé 520~millions~; la maîtrise d'ouvrage déclare 50~\% d'avancement,
mais seuls 40~\% des livrables sont recettés.
\begin{enumerate}[(a)]
    \item Calculez VA, IPC et EAC avec l'avancement déclaré.
    \item Recommencez avec l'avancement adossé aux livrables recettés.
    \item Quel chiffre présentez-vous au conseil, et sous quelle forme~?
    \item Quelle cotation retenir pour le contrôle « le suivi financier du
          projet est fiable »~?
\end{enumerate}
\end{boxexo}

\begin{boxexo}[title={\ding{45}\quad Exercice~19.3. L'audit de jalon d'HORIZON \niveauexpert}]
Placez-vous le 21~février~2026, au lendemain de la note de la DSI.
\begin{enumerate}[(a)]
    \item Quels faits un audit de jalon aurait-il établis, et avec quelles
          pièces~?
    \item Rédigez l'avis écrit de l'auditrice au comité d'audit, sans
          prendre la décision de bascule.
    \item Quelles cotations du Tableau de la lecture « Projet HORIZON »
          auraient pu changer si cet avis avait été suivi~?
    \item En quoi cet avis aurait-il réduit l'incertitude de la direction
          générale au moment de décider~?
\end{enumerate}
\end{boxexo}

\bigskip

\begin{boxexo}[title={\ding{45}\quad Exercice~19.4. Signaux de dérive \niveaubase}]
Citez quatre signaux de dérive d'un projet, et pour chacun la pièce qui
permet de le vérifier.
\end{boxexo}

\begin{boxexo}[title={\ding{45}\quad Exercice~19.5. La décision de mise en production \niveaumoyen}]
La bascule du système central de la BANQUE SANAGA a été décidée par
courriel du DSI, sans comité. Rédigez la fiche d'assertion du contrôle
« la mise en production est décidée par l'instance prévue, au vu de la
recette » et cotez-la.
\end{boxexo}

\begin{boxexo}[title={\ding{45}\quad Exercice~19.6. Indice de performance \niveaumoyen}]
Un projet a un budget de 800~millions, une valeur planifiée à date de
400~millions, une valeur acquise de 300~millions et un coût réel de
450~millions. Calculez les indices de performance des coûts et des délais,
et l'estimation à terminaison. Concluez.
\end{boxexo}

\begin{boxexo}[title={\ding{45}\quad Exercice~19.7. Le cahier des charges \niveaumoyen}]
Réécrivez de façon neutre l'exigence « compatibilité native avec le
produit de l'éditeur Y », et dites ce que l'auditeur doit vérifier si
l'organisation maintient l'exigence.
\end{boxexo}

\begin{boxexo}[title={\ding{45}\quad Exercice~19.8. La réversibilité \niveaumoyen}]
Le contrat de l'ANPP ne prévoit ni remise des codes sources ni séquestre.
Quel est le risque, et quelle clause proposez-vous~?
\end{boxexo}

\begin{boxexo}[title={\ding{45}\quad Exercice~19.9. Offre unique \niveauexpert}]
Pour le marché de la dématérialisation, concevez les tests qui
permettraient de passer de \NEK{} à \OUI{} ou à \NON{} sur le contrôle
« la mise en concurrence a été effective ». Quel est le seuil de preuve
pour coter \NON{} sans qualifier d'infraction~?
\end{boxexo}

\begin{boxexo}[title={\ding{45}\quad Exercice~19.10. Arrêter ou continuer~? \niveauexpert}]
Avec un indice de performance des coûts de 0,55 sur l'ensemble des trois
projets de l'ANPP, que peut recommander la Chambre des comptes~?
Distinguez ce qui relève du constat de l'auditeur et ce qui relève de la
décision du ministère.
\end{boxexo}

\bigskip

\begin{boxinfo}
\textbf{Points de contrôle de l'élève.} Je sais~: décomposer l'audit d'un
projet~; auditer une décision de mise en production~; calculer la valeur
acquise et ses indices~; auditer un marché étape par étape~; repérer un
cahier des charges qui oriente la concurrence.
\end{boxinfo}

\begin{boxinfo}
\textbf{Pour aller plus loin.}
\begin{itemize}
    \item Le guide du CHAI~\cite{chai2014guide}~: les fiches « projets » et
          « marchés informatiques ».
    \item ISACA, COBIT~2019~\cite{isaca2018cobitobjectifs}~: les objectifs
          BAI (bâtir, acquérir, implémenter).
\end{itemize}
\end{boxinfo}

\begin{boxinfo}
\textbf{Transition.} Une fois en production, une application repose sur une
infrastructure et sur des mesures de sécurité. Le
Chapitre~\ref{chap:audit-technique} audite la sécurité et l'infrastructure,
dont l'audit de sécurité légal et l'audit d'un SMSI.
\end{boxinfo}
